\section{总结与延伸}

\subsection{Visual Prompting 的统一范式}

本期介绍的所有工作——SoM、FreeGrasp、Pivot、VLM+SAM Agent——都遵循一个统一的范式：

\begin{importantbox}{Visual Prompting 统一范式}
\begin{enumerate}
    \item 使用经典 CV 模块（分割、检测、采样等）对原始图片做预处理
    \item 在预处理结果上叠加视觉标记（编号、箭头、掩码等）
    \item 将增强后的图片交给 VLM 做高层推理与决策
\end{enumerate}
这种"CV 预处理 + VLM 推理"的组合，充分利用了两者的互补优势。
\end{importantbox}

\subsection{技术演化方向}

\begin{itemize}
    \item \textbf{模型能力增强}：随着 VLM 视觉理解能力的提升，部分 Visual Prompting 场景可能不再必要，但在复杂场景中仍有价值
    \item \textbf{端到端整合}：未来可能出现将 CV 预处理与 VLM 推理端到端整合的统一模型
    \item \textbf{更多应用场景}：从物体定位、机器人抓取、导航，到自动驾驶、医学影像等更多领域
    \item \textbf{动态视频处理}：结合 SAM 2 的视频追踪能力，Visual Prompting 可以扩展到视频理解与交互
\end{itemize}

\subsection{拓展阅读}

\begin{itemize}
    \item Set-of-Mark Visual Prompting for GPT-4V（SoM 原论文）
    \item Segment Anything Model（SAM 1/2/3）系列论文
    \item Grounding DINO：开放词汇目标检测
    \item FreeGrasp：机器人自由形式抓取
    \item Pivot：迭代式 Visual Prompting 导航
\end{itemize}

%% --- Fin del contenido --- %%

\end{document}
